\mathtgt
\section{Canonical Forms of Closed Evidence}\label{sec:canonical}

Over a typed store every binder is transparent, every block name is defined by a stored literal, and evidence has no assumptions beyond those definitions. This section shows that such evidence normalizes, and that its normal forms have the shapes that preservation and progress need.

\subsection{Shapes and resolution}\label{sec:resolution}

A head form promises a shape for its endpoints, which is $\bot$, $\top$, a function type, or an object type, and shapes are read after following the definitions of the store. Write $\Gamma(x) \ni \ell \mathrel{:=} W$ when $x$ is transparent in $\Gamma$ with witness $W$ at $\ell$. One step of resolution, $\mathit{next}_\Gamma(T)$, is $W$ when $T$ is a name $\sel{x}{\ell}$ with $\Gamma(x) \ni \ell \mathrel{:=} W$, and is undefined when $T$ is any other type or the name of an opaque binder. A chain of definitions may be cyclic, because a witness may itself be a name of the same block, as in the literal
\[
\lit{A = \sel{\self}{B},\ B = \sel{\self}{A}}{\tnil} .
\]
Resolution therefore runs with fuel:
\[
\begin{array}{@{}lcl@{}}
\resolven{\Gamma}{0}{T} & = & \top \text{ if } \mathit{next}_\Gamma(T) \text{ defined, else } T \\
\resolven{\Gamma}{n+1}{T} & = & \resolven{\Gamma}{n}{W} \text{ if } \mathit{next}_\Gamma(T) = W, \text{ else } T \\
\resolve{\Gamma}{T} & = & \resolven{\Gamma}{|\mathit{defs}(\Gamma)| + 1}{T}
\end{array}
\]
where $\mathit{defs}(\Gamma)$ lists the pairs of a transparent binder and a label its witness list mentions. A name at a label the list does not mention is defined as $\top$ and settles in one step. A chain that runs out of fuel on a defined name is cyclic, and a cyclic alias resolves to $\top$, the object type that asserts nothing. The only facts known about the names of a cycle are their equalities with each other, and $\top$ asserts nothing more. Resolving to $\bot$ instead would be unsound: a field of the literal above declared at $\sel{\self}{A}$ holds a value, an atom bound to that field has a type that resolves through the cycle, and item 5 of Theorem~\ref{thm:canon} says that the type of a well-typed atom never resolves to $\bot$. The fuel is enough: a chain of definitions either settles, at a shape or at an undefined name, within $|\mathit{defs}(\Gamma)| + 1$ steps, or it visits a defined name twice and is periodic from then on. This is a pigeonhole argument on the finitely many defined names, and it gives the two facts used everywhere below, with no side condition on the context.

\begin{lemma}[Resolution]\label{lem:resolve}
\begin{enumerate}
\item If $\Gamma(x) \ni \ell \mathrel{:=} W$ then
\[ \resolve{\Gamma}{\sel{x}{\ell}} = \resolve{\Gamma}{W} . \]
\item For every $T$,
\[ \resolve{\Gamma}{\resolve{\Gamma}{T}} = \resolve{\Gamma}{T} . \]
\end{enumerate}
\end{lemma}

The first version of the development required definitions to be well founded, with a side condition on literals that a witness may not be a bare name of the same block, the self-alias restriction. The condition propagated into the translation and, from there, into a restriction on the source rule \rn{\{\}-I}. Replacing well-founded resolution by the fuel-and-pigeonhole one removed the condition on both sides. Section~\ref{sec:forced} returns to this.

\subsection{Forms, entries, and views}

\begin{figure*}
\small
\color{tgtcol}
\[
\begin{array}{@{}lrcl@{}}
\text{Forms} & F,G,R & ::= & \form{bot} \mid \form{top} \mid \form{id} \mid \form{eqv}\,\varphi \mid \form{pi}\,d\,c \mid \form{obj}\,\mathit{Es} \mid \form{bnd}\,i\,F \mid \form{into}\,\mathit{Es} \\
\text{Entries} & E & ::= & \form{le}\,F\,H\,G \mid \form{eq}\,j\,b \mid \form{has}\,j \mid \form{bnd}\,G \mid \form{thru}\,R\,E \qquad \mathit{Es} ::= \tnil \mid \tcons{\mathit{Es}}{E} \\
\text{Proposition forms} & Q & ::= & \form{le}\,F \mid \form{eq} \mid \form{has}\,x\,\ell \mid \form{bnd}\,G \qquad\quad V ::= \tnil \mid \tcons{V}{Q}
\end{array}
\]
\color{black}
\caption{Head forms of closed inclusion evidence, their entries, and views of atoms.}
\label{fig:forms}
\end{figure*}

Figure~\ref{fig:forms} gives the normal forms: the head form of an inclusion, the entry of a template, and the view of an atom. Head forms wear a hat, to tell $\form{obj}$, the normal form, from $\ev{obj}$, the evidence. A head form of an inclusion is $\form{bot}$, $\form{top}$, the syntactic identity $\form{id}$, a conversion $\form{eqv}\,\varphi$ by closed equality evidence, a function coercion $\form{pi}\,d\,c$, or an object coercion $\form{obj}\,\mathit{Es}$ with one entry per proposition of the target telescope. The domain and codomain evidence of a function form is not normalized further, because the machine uses it as a cast and normalizes it only when it is applied. An entry $\form{le}\,F\,H\,G$ is the normal form of a template: the two sides are forms, $\form{id}$ standing for an absent side, around a hole $H$ that still names a source proposition. The other entries are as in morphisms.

Two head forms and two entries belong to self-bounds. The form $\form{bnd}\,i\,F$ casts through the $i$-th bound of the source object type and continues by $F$ from the type of that bound. The form $\form{into}\,\mathit{Es}$ is a coercion into an object type whose entries never consult the view of the source. Its entries are bound entries $\form{bnd}\,G$, closed coercions out of the source, and routed entries $\form{thru}\,R\,E$, in which the form $R$ reaches some object type from the source and the entry $E$ proves the target proposition there. Routes never nest, and only $\form{into}$ forms carry routed entries. A route is a subterm of the form that carries it, so applying a form to a view recurses on smaller forms.

A view is the telescope of forms of the propositions that an atom's resolved object type asserts about it. The form of an inclusion proposition is a head form, an equality proposition records only that its two sides resolve equally, and a presence proposition records the binder and the field, $\form{has}\,x\,\ell$. A bound proposition records a form $\form{bnd}\,G$ from the type of the root into the bound, since a bound of an opened telescope is a fact about the root.

\subsection{The normalizer}\label{sec:normalizer}

\begin{figure*}
\small
\color{tgtcol}
\[
\begin{array}{@{}lcl@{}}
\mathit{hnf}_\sigma(\ev{refl}\,T) & = & \form{eqv}\,(\ev{refl}\,T) \\
\mathit{hnf}_\sigma(\ev{top}\,T) & = & \form{top} \qquad\qquad \mathit{hnf}_\sigma(\ev{bot}\,T) \ =\ \form{bot} \\
\mathit{hnf}_\sigma(\ev{eqToLe}\,\varphi) & = & \form{eqv}\,\varphi \qquad\quad\ \mathit{hnf}_\sigma(\ev{pi}\,d\,c) \ =\ \form{pi}\,d\,c \\
\mathit{hnf}_\sigma(\ev{obj}\,\tel\,m) & = & \form{obj}\,(\mathit{entries}_\sigma(m)) \\
\mathit{hnf}_\sigma(\ev{pair}\,\tel_1\,\tel_2\,e\,f) & = & \mathit{pair}_{\tel_1,\tel_2}\,(\mathit{hnf}_\sigma(e))\,(\mathit{hnf}_\sigma(f)) \\
\mathit{hnf}_\sigma(\ev{trans}\,e\,f) & = & \mathit{hnf}_\sigma(e) \circ \mathit{hnf}_\sigma(f) \\
\mathit{hnf}_\sigma(\ev{bound}\,\tel\,i) & = & \form{bnd}\,i\,\form{id} \qquad\quad \mathit{hnf}_\sigma(\ev{intoBnd}\,e) \ =\ \form{into}\,[\,\form{bnd}\,(\mathit{hnf}_\sigma(e))\,] \\
\mathit{hnf}_\sigma(\ev{member}\,a\,e\,i) & = & G \quad \text{where } \mathit{viewThrough}_\sigma(\mathit{hnf}_\sigma(e), a) \ni (i \mapsto \form{le}\,G) \\[3pt]
\mathit{view}_\sigma(x) & = & \text{the precise view of } \sigma(x) \text{ at } x \\
\mathit{view}_\sigma(\ev{cast}\,a\,e) & = & \mathit{viewThrough}_\sigma(\mathit{hnf}_\sigma(e), a) \\
\mathit{view}_\sigma(\ev{fold}\,\tel\,a) & = & \mathit{view}_\sigma(a) \ =\ \mathit{view}_\sigma(\ev{unfold}\,a) \\
\mathit{view}_\sigma(\ev{both}\,\tel_1\,\tel_2\,a\,b) & = & \mathit{view}_\sigma(a) \mathbin{+\!\!+} \mathit{view}_\sigma(b) \\[3pt]
\mathit{viewThrough}_\sigma(\form{id}, a) & = & \mathit{view}_\sigma(a) \ =\ \mathit{viewThrough}_\sigma(\form{eqv}\,\varphi, a) \\
\mathit{viewThrough}_\sigma(\form{obj}\,\mathit{Es}, a) & = & \mathit{entriesAt}_\sigma(a, \mathit{Es}) \ =\ \mathit{viewThrough}_\sigma(\form{into}\,\mathit{Es}, a) \\
\mathit{viewThrough}_\sigma(\form{bnd}\,i\,F, a) & = & \mathit{viewThrough}_\sigma(G \circ F, \rt{a}) \quad \text{where } \mathit{view}_\sigma(a) \ni (i \mapsto \form{bnd}\,G) \\
\mathit{viewThrough}_\sigma(F, a) & = & \tnil \quad \text{for } F \in \{\form{pi}\,d\,c, \form{top}, \form{bot}\}
\end{array}
\]
\color{black}
\caption{The normalizer, with the fuel argument and the failure cases elided. Each recursive call is on a subterm and consumes one unit of fuel. $\mathit{entriesAt}_\sigma(a, \mathit{Es})$ instantiates the entries at the view of $a$ and, for bound and routed entries, at the chain of casts of $a$.}
\label{fig:normalizer}
\end{figure*}

\begin{figure*}
\small
\color{tgtcol}
\[
\begin{array}{@{}lcl@{}}
\form{id} \circ G = G & \qquad & F \circ \form{id} = F \\
\form{bot} \circ G = \form{bot} & & F \circ \form{top} = \form{top} \\
\form{eqv}\,\varphi \circ \form{eqv}\,\psi = \form{eqv}\,(\ev{trans}\,\varphi\,\psi) & & \form{eqv}\,\varphi \circ F = F = F \circ \form{eqv}\,\varphi \quad \text{for } F \in \{\form{pi}\,d\,c, \form{obj}\,\mathit{Es}\} \\
\multicolumn{3}{@{}l@{}}{\form{pi}\,d_1\,c_1 \circ \form{pi}\,d_2\,c_2 = \form{pi}\,(\ev{trans}\,d_2\,d_1)\,(\ev{trans}\,(c_1[x \mapsto \ev{cast}\,x\,d_2])\,c_2)} \\
\multicolumn{3}{@{}l@{}}{\form{obj}\,\mathit{Es}_1 \circ \form{obj}\,\mathit{Es}_2 = \form{obj}\,(\mathit{Es}_2 \text{ routed through } \mathit{Es}_1)} \\
\multicolumn{3}{@{}l@{}}{\form{bnd}\,i\,F \circ G = \form{bnd}\,i\,(F \circ G) \qquad\qquad \form{obj}\,\mathit{Es} \circ \form{bnd}\,i\,F = G \circ F \quad \text{where } \mathit{Es} \ni (i \mapsto \form{bnd}\,G)} \\
\multicolumn{3}{@{}l@{}}{F \circ \form{into}\,\mathit{Es} = \form{into}\,(\mathit{Es} \text{ prefixed by } F) \qquad \form{into}\,\mathit{Es}_1 \circ \form{obj}\,\mathit{Es}_2 = \form{into}\,(\mathit{Es}_2 \text{ prefixed by } \form{into}\,\mathit{Es}_1) \qquad \form{top} \circ \form{obj}\,\mathit{Es} = \form{into}\,(\mathit{Es} \text{ prefixed by } \form{top})}
\end{array}
\]
\color{black}
\caption{Composition of head forms, $F \circ G$. In addition, $\form{eqv}\,\varphi \circ \form{bot} = \form{bot}$, conversions are absorbed by $\form{bnd}$ and $\form{into}$ forms as they are by $\form{pi}$ and $\form{obj}$ forms, an $\form{into}$ form composed with a bound cast reads its entry at the bound like an $\form{obj}$ form, and in every other case $F \circ G = F$. Routing the entries of one object form through another fails when an entry names a proposition of the wrong kind, which does not happen between typed forms.}
\label{fig:combine}
\end{figure*}

Figure~\ref{fig:normalizer} gives the normalizer, a family of mutually recursive functions indexed by fuel. Each returns an option, and the fuel bound of a successful run is determined by the size of the input, so the functions are total and the fuel is bookkeeping. Normalization of an inclusion is written $\hnf{\sigma}{e}{n}{F}$, the view of an atom $\vw{\sigma}{a}{n}{V}$, and the chain of casts of an atom $\chn{\sigma}{a}{n}{(a', F)}$. The results are monotone in the fuel and deterministic.

The precise view of a stored literal lists $\form{eq}$ for each witness and $\form{has}\,x\,\ell$ for each field, in the order of the literal's precise telescope. The view of a cast atom is the view of the underlying atom pushed through the head form of the cast, and folding, unfolding, and pairing act on views as they act on telescopes.

Composition, $F \circ G$ in Figure~\ref{fig:combine}, is where the design of Section~\ref{sec:why-structural} shows. Conversions are absorbed. Two function forms compose contravariantly in the domain and by substituting the domain cast for the parameter in the first codomain. Two object forms compose by routing every entry of the second through the entries of the first. Consider an entry of the second form whose hole names the proposition $j$ of the middle telescope. The first form's entry at position $j$ is itself a template around a hole in the source telescope, or it is an equality. Either way the composite entry keeps that hole and composes the sides. Every recursive call to composition is on a strict subterm, and no view is consulted.

Prefixing an entry by a form $F$ makes it read the object type that $F$ reaches. A bound entry $\form{bnd}\,G$ becomes $\form{bnd}\,(F \circ G)$, a routed entry $\form{thru}\,R\,E$ becomes $\form{thru}\,(F \circ R)\,E$, and any other entry becomes $\form{thru}\,F\,E$. The one exception is the identity on a source bound, $\form{bnd}\,(\form{bnd}\,j\,\form{id})$, which is kept routed rather than composed, so that a later bound cast through position $j$ still finds the bound it names. A bound cast composed with an object form reads the entry at its index and continues from there.

Application of an object form to a view, $\mathit{entriesAt}$, is pointwise. A template is instantiated by replacing its hole by the view's form of the proposition it names, an equality reading as $\form{id}$, and composing with the sides. An equality entry requires the view to hold an equality at the named position, and a presence entry copies the binder and label found there. A bound entry $\form{bnd}\,G$ is instantiated by composing the chain of casts of the atom with $G$, which is why $\mathit{entriesAt}$ takes the atom and not only its view. A routed entry $\form{thru}\,R\,E$ recomputes the view of the atom through $R$ and instantiates $E$ there.

Pairing produces an $\form{into}$ form. When either component is $\form{bot}$ the pair is $\form{bot}$, when both are $\form{top}$ it is $\form{top}$, and a component that is $\form{bot}$ behind bound casts is typed evidence for every inclusion out of its source and is returned as it is. Otherwise each component contributes its entries, routed through the identity or through the bound cast the component goes under, and the pair concatenates them. A component may be a bound cast $\form{bnd}\,i\,F$, whose target propositions are reached through an object type other than the source, and this is the case the routed entries exist for. An $\form{obj}$ form names source propositions by index and cannot express it.

The chain of casts of an atom, $\chn{\sigma}{a}{n}{(a', F)}$, composes the head forms of the casts on the atom from its root outwards, pairs at $\ev{both}$, and passes through $\ev{fold}$ and $\ev{unfold}$. It is the form the machine consults at an application.

\subsection{Typed forms}

\begin{figure*}
\small
\color{tgtcol}
\begin{mathpar}
\inferrule{\resolveAt{\Gamma}{\rho}{S} = \bot}{\chainj{\Gamma}{\rho}{\form{bot}}{S}{T}}
\and
\inferrule{\resolveAt{\Gamma}{\rho}{T} = \top}{\chainj{\Gamma}{\rho}{\form{top}}{S}{T}}
\and
\inferrule{\resolveAt{\Gamma}{\rho}{S} = \resolveAt{\Gamma}{\rho}{T}}{\chainj{\Gamma}{\rho}{\form{id}}{S}{T}}
\and
\inferrule{\resolveAt{\Gamma}{\rho}{S} = \resolveAt{\Gamma}{\rho}{T}}{\chainj{\Gamma}{\rho}{\form{eqv}\,\varphi}{S}{T}}
\and
\inferrule{\resolveAt{\Gamma}{\rho}{S} = \pit{S_1}{T_1} \\ \resolveAt{\Gamma}{\rho}{T} = \pit{S_2}{T_2} \\ \lej{\Gamma}{d}{S_2}{S_1} \\ \lej{\Gamma, x : S_2}{c}{T_1}{T_2}}
  {\chainj{\Gamma}{\rho}{\form{pi}\,d\,c}{S}{T}}
\and
\inferrule{\resolveAt{\Gamma}{\rho}{S} = \obj{\tel_1} \\ \resolveAt{\Gamma}{\rho}{T} = \obj{\tel_2} \\ \entjm{\Gamma}{\rho}{\mathit{Es}}{\tel_1}{\tel_2}}
  {\chainj{\Gamma}{\rho}{\form{obj}\,\mathit{Es}}{S}{T}}
\and
\inferrule{\resolveAt{\Gamma}{\rho}{S} = \obj{\tel} \\ \tat{\tel}{i}{\pbnd{\weak{T}}} \\ \chainj{\Gamma}{\rho}{F}{T}{U}}
  {\chainj{\Gamma}{\rho}{\form{bnd}\,i\,F}{S}{U}}
\and
\inferrule{\resolveAt{\Gamma}{\rho}{T} = \obj{\tel} \\ \entjm{\Gamma}{\rho}{\mathit{Es}}{S}{\tel}}
  {\chainj{\Gamma}{\rho}{\form{into}\,\mathit{Es}}{S}{T}}
\\
\inferrule{ }{\entjm{\Gamma}{\rho}{\tnil}{\tel_1}{\tnil}}
\and
\inferrule{\entjm{\Gamma}{\rho}{\mathit{Es}}{\tel_1}{\tel_2} \\ \holej{\tel_1}{H}{X}{Y} \\ \Gamma \models F : S \le X \text{ (side)} \\ \Gamma \models G : Y \le T \text{ (side)}}
  {\entjm{\Gamma}{\rho}{\tcons{\mathit{Es}}{\form{le}\,F\,H\,G}}{\tel_1}{\tcons{\tel_2}{\ple{S}{T}}}}
\and
\inferrule{\entjm{\Gamma}{\rho}{\mathit{Es}}{\tel_1}{\tel_2} \\ \tat{\tel_1}{j}{\phas{\ell}}}
  {\entjm{\Gamma}{\rho}{\tcons{\mathit{Es}}{\form{has}\,j}}{\tel_1}{\tcons{\tel_2}{\phas{\ell}}}}
\and
\inferrule{\entjm{\Gamma}{\rho}{\mathit{Es}}{\tel_1}{\tel_2} \\ \chainj{\Gamma}{\rho}{G}{\obj{\tel_1}}{T}}
  {\entjm{\Gamma}{\rho}{\tcons{\mathit{Es}}{\form{bnd}\,G}}{\tel_1}{\tcons{\tel_2}{\pbnd{\weak{T}}}}}
\and
\inferrule{\entjm{\Gamma}{\rho}{\mathit{Es}}{\tel_1}{\tel_2} \\ \tat{\tel_1}{j}{\pbnd{X}}}
  {\entjm{\Gamma}{\rho}{\tcons{\mathit{Es}}{\form{bnd}\,(\form{bnd}\,j\,\form{id})}}{\tel_1}{\tcons{\tel_2}{\pbnd{X}}}}
\\
\inferrule{ }{\entjm{\Gamma}{\rho}{\tnil}{S}{\tnil}}
\and
\inferrule{\entjm{\Gamma}{\rho}{\mathit{Es}}{S}{\tel} \\ \chainj{\Gamma}{\rho}{F}{S}{T}}
  {\entjm{\Gamma}{\rho}{\tcons{\mathit{Es}}{\form{bnd}\,F}}{S}{\tcons{\tel}{\pbnd{\weak{T}}}}}
\and
\inferrule{\entjm{\Gamma}{\rho}{\mathit{Es}}{S}{\tel} \\ \chainj{\Gamma}{\rho}{R}{S}{M} \\\\ \resolveAt{\Gamma}{\rho}{M} = \obj{\tel_M} \\ \entjm{\Gamma}{\rho}{E}{\tel_M}{P}}
  {\entjm{\Gamma}{\rho}{\tcons{\mathit{Es}}{\form{thru}\,R\,E}}{S}{\tcons{\tel}{P}}}
\\
\inferrule{ }{\viewj{\Gamma}{r}{\sigma}{\tnil}{\tnil}}
\and
\inferrule{\viewj{\Gamma}{r}{\sigma}{V}{\tel} \\ \formj{\Gamma}{F}{\inst{S}{r}}{\inst{T}{r}}}
  {\viewj{\Gamma}{r}{\sigma}{\tcons{V}{\form{le}\,F}}{\tcons{\tel}{\ple{S}{T}}}}
\and
\inferrule{\viewj{\Gamma}{r}{\sigma}{V}{\tel} \\ \resolve{\Gamma}{\inst{S}{r}} = \resolve{\Gamma}{\inst{T}{r}}}
  {\viewj{\Gamma}{r}{\sigma}{\tcons{V}{\form{eq}}}{\tcons{\tel}{\peq{S}{T}}}}
\and
\inferrule{\viewj{\Gamma}{r}{\sigma}{V}{\tel} \\ \sigma(r) \text{ has field } \ell}
  {\viewj{\Gamma}{r}{\sigma}{\tcons{V}{\form{has}\,r\,\ell}}{\tcons{\tel}{\phas{\ell}}}}
\and
\inferrule{\viewj{\Gamma}{r}{\sigma}{V}{\tel} \\ \chainj{\Gamma}{r}{G}{\Gamma(r)}{\inst{X}{r}}}
  {\viewj{\Gamma}{r}{\sigma}{\tcons{V}{\form{bnd}\,G}}{\tcons{\tel}{\pbnd{X}}}}
\end{mathpar}
\color{black}
\caption{Typed forms $\chainj{\Gamma}{\rho}{F}{S}{T}$ in a mode $\rho$, typed entries $\entjm{\Gamma}{\rho}{\mathit{Es}}{\tel_1}{\tel_2}$, and typed views $\viewj{\Gamma}{r}{\sigma}{V}{\tel}$. The entry rules for equalities mirror the morphism rules. A side is typed as $\form{id}$ between equal types, or as a form in the plain mode between closed types weakened under the self binder. The entries of an $\form{into}$ form are typed from the source type rather than from its telescope, since they never consult the view, and $\entjm{\Gamma}{\rho}{E}{\tel_M}{P}$ types one entry at one proposition.}
\label{fig:formtyped}
\end{figure*}

A head form is typed syntactically, in Figure~\ref{fig:formtyped}: it is typed evidence for $S \le T$ when its pieces of evidence are typed and the two endpoints resolve to the shapes the form promises. The entries of an object form are typed between two closed telescopes, entry by entry, against the morphism rules of Figure~\ref{fig:evidence}. The domain and codomain of a function form are typed as evidence, since that is how the machine uses them. A view is typed against a telescope instantiated at the root of the atom.

Typedness has two modes, indexed by an optional root $\rho$, and keeping them apart is what makes the elimination case of the theorem go through. In the plain mode, written $\formj{\Gamma}{F}{S}{T}$, the shape of a type is its resolution, and a coercion form is typed between closed types unrelated to any particular atom. This is what lets the $\ev{member}$ case of the theorem extract a form from a view and reuse it as a coercion at no root at all. In the mode at a root $r$ the shape of a type is its resolution with the self block opened at that root:
\[
\begin{array}{@{}rcl@{}}
\resolveAt{\Gamma}{r}{T} & = & \mathit{open}_r(\resolve{\Gamma}{T}) \\
\mathit{open}_r(\obj{\tel}) & = & \obj{\weak{(\inst{\tel}{r})}} \qquad \mathit{open}_r(T) = T \text{ otherwise} .
\end{array}
\]
The chain of casts of an atom is typed in the mode at its root, and opening makes $\ev{fold}$ and $\ev{unfold}$ invisible to it. The two modes are one inductive family, and every lemma about typed forms is proven once for both. The only rule that mixes them is the view rule for a bound: the form of a bound in the view of an atom is typed at the root, from the type of the root into the bound, because a bound of an opened telescope is a fact about the root. The first version of the development defined the rooted judgment as the plain one at opened endpoints. Bounds ended that, because the form that continues from a bound has the bound's type as its source, and that type has to be opened at the root as well.

\subsection{Composition, pairing, and application preserve typedness}

\begin{lemma}[Composition]\label{lem:combine}
If, in one mode $\rho$,
\[ \chainj{\Gamma}{\rho}{F}{S}{M} \quad\text{and}\quad \chainj{\Gamma}{\rho}{G}{M}{T} \]
then $F \circ G$ is defined and $\chainj{\Gamma}{\rho}{F \circ G}{S}{T}$.
\end{lemma}

The proof is an induction on the sum of the sizes of the two forms, mutually with the analogous statement for routing typed entries through typed entries and for prefixing typed entries by a typed form. The object case is the one that matters. A template of the second form has a hole naming a proposition of the middle telescope, and the first form's typed entry at that position is a template around a hole in the source telescope, or an equality. The sides compose by the induction hypothesis at smaller size, and because sides are closed forms typed between weakened closed types, the composite never mentions the self block. Two function forms compose because the domain casts compose in the reverse direction and the substitution of a typed cast for the parameter preserves the typing of the codomain evidence. A bound cast composed with a form continues by the composite, and an object form composed with a bound cast reads its entry at the bound, which is a strict subterm.

\begin{lemma}[Pairing and application]\label{lem:pair-apply}
If $\chainj{\Gamma}{\rho}{F}{S}{\obj{\tel_1}}$ and $\chainj{\Gamma}{\rho}{G}{S}{\obj{\tel_2}}$ then $\mathit{pair}_{\tel_1,\tel_2}\,F\,G$ is defined and typed at
\[ S \le \obj{(\tel_1 \mathbin{+\!\!+} \tel_2)} \]
in the same mode. If $\entjm{\Gamma}{r}{\mathit{Es}}{\tel_1}{\tel_2}$, $\viewj{\Gamma}{r}{\sigma}{V}{\tel_1}$, and the chain of casts of $a$ is typed at $r$ from $\Gamma(r)$ to a type that opens to $\obj{\tel_1}$, then $\mathit{entriesAt}_\sigma(a, \mathit{Es})$ is defined and is a view typed at $\tel_2$.
\end{lemma}

The pairing lemma reads the entries of each component free of the source's view, which is possible because an $\form{obj}$ form's entries can be routed through the identity and a bound cast's entries through the bound.

\subsection{The theorem}

\begin{theorem}[Canonical forms]\label{thm:canon}
Let $\storej{\sigma}{\Gamma}$.
\begin{enumerate}
\item If $\lej{\Gamma}{e}{S}{T}$ then $\hnf{\sigma}{e}{n}{F}$ for some $n$ and $F$ with $\formj{\Gamma}{F}{S}{T}$.
\item If $\eqj{\Gamma}{\varphi}{S}{T}$ then $\resolve{\Gamma}{S} = \resolve{\Gamma}{T}$.
\item If $\hasj{\Gamma}{h}{x}{\ell}$ then presence normalization of $h$ at $x$ succeeds and $\sigma(x)$ is a literal with a field $\ell$.
\item If $\morj{\Gamma}{m}{\tel}{\tel'}$ then the entries of $m$ normalize to $\mathit{Es}$ with $\entj{\Gamma}{\mathit{Es}}{\tel}{\tel'}$.
\item If $\atomj{\Gamma}{a}{S}$ then $\vw{\sigma}{a}{n}{V}$ for some $n$ and $V$, with $\resolve{\Gamma}{S} \neq \bot$, and
\[ \viewj{\Gamma}{\rt{a}}{\sigma}{V}{\tel} \]
whenever $\resolve{\Gamma}{S} = \obj{\tel}$.
\end{enumerate}
\end{theorem}

\begin{proof}
By mutual structural induction on the evidence derivations, together with Theorem~\ref{thm:chain} below, because the view of an atom through a bound entry reads the chain of casts of the atom. The cases for $\ev{refl}$, $\ev{top}$, $\ev{bot}$, $\ev{eqToLe}$, and $\ev{pi}$ produce their form with one unit of fuel. The case $\ev{bound}$ produces $\form{bnd}\,i\,\form{id}$, and $\ev{intoBnd}$ produces $\form{into}\,[\,\form{bnd}\,F\,]$ from the form $F$ of its premise. The case $\ev{trans}$ is Lemma~\ref{lem:combine}, $\ev{pair}$ is Lemma~\ref{lem:pair-apply}, and $\ev{obj}$ is item 4 applied to the morphism, whose $\ev{bnd}$ entries become bound entries by item 1. For $\ev{member}\,a\,e\,i$, item 5 gives a typed view of $a$, and item 1 gives a typed form of $e$ whose target is an object type. Hence the source of $e$ resolves to an object type, and the view of $a$ through the form of $e$ is a typed view of the target telescope by Lemma~\ref{lem:pair-apply}. Its entry at $i$ is a typed form, an equality of resolutions, or a field of the object at the root, according to the kind of the proposition. The case $\ev{def}$ is Lemma~\ref{lem:resolve}. For an atom that is a variable, the precise view of the stored literal is typed at the literal's telescope, because each definition equality resolves equally by Lemma~\ref{lem:resolve} and each presence names a field the literal has. A cast atom pushes its view through the form of the cast, and through a bound cast it is the view of the root through the composite of the bound's form and the rest of the cast. Folding and unfolding keep the view, since a typed view of $\tel$ is a typed view of the opened telescope and conversely. The fuel of every case is at most two more than the maximum of the fuels of the premises.
\end{proof}

\begin{theorem}[Chain of casts]\label{thm:chain}
If $\storej{\sigma}{\Gamma}$ and $\atomj{\Gamma}{a}{S}$ then $\chn{\sigma}{a}{n}{(a', F)}$ for some $n$, $a'$, and $F$ with
\[ \chainj{\Gamma}{\rt{a}}{F}{\Gamma(\rt{a})}{S} . \]
\end{theorem}

The chain theorem types the composite of the casts of an atom from the type of its root to the type of the atom, at the root. Its proof composes the forms of the casts by Lemma~\ref{lem:combine}, lifted to the rooted judgment, and uses that opening at the root is stable under folding and unfolding.

\begin{corollary}[Shapes of closed inclusions]\label{cor:shapes}
Let $\storej{\sigma}{\Gamma}$ and $\lej{\Gamma}{e}{S}{T}$. Then $S$ resolves to $\bot$, or $T$ resolves to $\top$, or $S$ and $T$ resolve equally, or both resolve to function types, or both resolve to object types, or $S$ resolves to an object type one of whose bounds is included in $T$ by a typed form, or $T$ resolves to an object type all of whose propositions are bounds. In particular there is no closed evidence for $\top \le \bot$, none from an object type without bounds into a function type, and none from a function type into an object type with a proposition that is not a bound.
\end{corollary}

\begin{corollary}[Inversion at the store]\label{cor:inversion}
Let $\storej{\sigma}{\Gamma}$.
\begin{enumerate}
\item If $\atomj{\Gamma}{a}{\pit{S}{T}}$, then $\sigma(\rt{a})$ is a lambda, and the chain of casts of $a$ normalizes to $\form{id}$, a conversion, or a function form.
\item If $\hasj{\Gamma}{h}{x}{\ell}$, then $\sigma(x)$ is a literal with a field $\ell$.
\item For every binder $x$ and label $\ell$ there is a $W$ with $\Gamma(x) \ni \ell \mathrel{:=} W$ and
\[ \eqj{\Gamma}{\ev{def}\,x\,\ell}{\sel{x}{\ell}}{W} . \]
\end{enumerate}
\end{corollary}

The first corollary is the target-side residue of what invertible typing provides in the DOT proofs, proven once, over stores, with no restriction on recursive members and no semantic model. The second is exactly what progress needs.

\subsection{Preservation, progress, and simulation}

\begin{theorem}[Preservation]
If a state is typed at $U$ and steps to a state over a possibly larger signature, the new state is typed at $U$ renamed into that signature.
\end{theorem}

The only rule whose case needs the store is application through a wrapped atom. There the machine reads a function form $\form{pi}\,d\,c$ off the chain of casts of the atom, and Theorem~\ref{thm:chain} types that form from the root's type to the atom's function type, at the root. The store holds a lambda at the root, so the root's type is syntactically a function type, and the shapes of the typed form give $d$ and $c$ the endpoints that the cast on the argument and the cast on the result need. Allocation uses that stripping the casts of a typed value leaves a typed literal together with a typed composite cast, and that a term typed with a binder opaque is typed with it transparent. The remaining cases are the usual renaming and substitution lemmas.

\begin{theorem}[Progress]
A typed state is final or steps.
\end{theorem}

The two rules that consult the store are application and projection, and Corollary~\ref{cor:inversion} supplies a closure at the root of a function atom with a normalizing chain of casts, and a field at the root of a projection.

\begin{theorem}[Erasure simulation]
A step of the machine either moves a cast frame and leaves the erasure unchanged, or erases to one runtime step. Conversely, if $\storej{\sigma}{\Gamma}$, the term of the state is typed, and the erased state takes a runtime step to $r$, then the state takes a run to a state whose erasure is $r$.
\end{theorem}

The backward direction realizes a runtime application by first moving the pending cast frames, then applying through the atom, which Corollary~\ref{cor:inversion} makes possible. Along any run from a typed state the store stays typed, so Corollaries~\ref{cor:shapes} and~\ref{cor:inversion} hold at every reachable state: no reachable store admits closed evidence for $\top \le \bot$, and every block name of every reachable store is defined.
